\section{BAO}\label{sec:bao}

\subsection{Data}
We use the 13 DESI DR2 BAO distance measurements of Results~II, Table~IV \citep{DESI2025a,DESI2025b}, spanning $z=0.295$--$2.33$ ($D_M/r_d$, $D_H/r_d$, $D_V/r_d$), with the full $13\times13$ covariance as distributed by the CobayaSampler \code{desi\_bao\_dr2} package. Here $D_H=c/H$ and the DESI convention is $D_V=(zD_M^2D_H)^{1/3}$; the general volume-averaged-distance construction is reviewed in \citet{Eisenstein2005}, but the prefactor and exponent placement used here are explicitly the DESI one. The scale $f=c/(H_0r_d)$ is profiled by Brent minimisation on the interval $[20,80]$. These reduced three-dimensional BAO distances inherit the fiducial-cosmology assumption of the data reduction. This model dependence, discussed for bimetric gravity by \citet{DwivediHogas2024}, is not removed by our likelihood; a 2D-BAO reanalysis is outside the present scope.

\subsection{Results}\label{sec:baores}
The bimetric model yields $\chi^2_{\rm BAO}=10.458$ against $10.279$ for \LCDM{}, i.e.\ $\Dchi=0.179$. The formal one-parameter tail value $p=0.672$ is retained only as a descriptive reference, not as a model-selection significance; Wilks-regularity conditions are not automatic here (in particular, at $\theta=0$ the parameter $\mFP$ becomes unidentifiable). At the adopted anchor, BAO is insensitive to the four varied bimetric parameters at current data precision:
\begin{itemize}
\item \emph{$(\alpha,\beta)$ scan} (19 non-degenerate combinations): $\Dchi=0.1791$--$0.1794$, flat across the entire grid. The degenerate locus $B_1=0$, $1+2\alpha+\beta=0$ (e.g.\ $(\alpha,\beta)=(-1,1)$), carries no finite branch and is reported as such; it is excluded from the scan rather than silently replaced by a GR value.
\item \emph{$\mFP/H_0$ scan} ($15$ to $10^5$, 3.8 decades): $\Dchi$ flat at $0.179$.
\item \emph{$\theta$ scan:} with $\Om$ refit at each angle (the self-consistent scan), $\Dchi$ decreases mildly from $0.179$ at $6.1\deg$ to $0.177$ at $31.5\deg$; profiling of $f$ absorbs the $\theta$-dependence even at large angles. Only a fixed-$\Om$ variant ($\Om$ frozen at its benchmark value) shows growth, reaching $\Dchi=55.3$ at $31.5\deg$; this variant overstates the sensitivity, since $\Om$ in the model is set by the quartic closure and is not held fixed. The benchmark $\theta=6.1\deg$ is deep in the blind regime.
\item \emph{Synthetic null.} Under the \LCDM{} null (2000 synthetic DESI DR2 realisations drawn from the full covariance, with $f$ refit for both models \emph{and} $\Om$ refit for \LCDM{}, matching the observed-data comparison), $\Dchi\ge0.179$ occurs with probability $p=0.682$: the BAO penalty is statistically unremarkable.
\end{itemize}

\paragraph{Mechanism: the dark degeneracy.}
The decomposition of $\Dchi$ ($+0.665$ from the density shift, $-0.486$ from the $E^2(z)$ difference, net $0.179$) is the bookkeeping of a near-identity. The late-time bimetric background is \LCDM{} with the effective density $\Om^{\rm eff}=\cos^2\!\theta\,\Om=0.30096$ (the dark degeneracy, \citealp{Kunz2009}, in the late-time effective form noted by \citealp{Hogas2021}), with a computed maximum absolute residual of $3.3288\times10^{-5}$ in $E^2$ at $z\le2.33$ (conservatively quoted elsewhere as $3.4\times10^{-5}$), negative throughout the shaded window (\cref{fig:mechanism}a) and a pure-correction penalty of $0.022$ in $\chi^2$ against that effective model. The BAO penalty is therefore a density-shift effect, not a shape-discrimination effect.

Two ``pure correction'' conventions exist in the literature, and we disclose ours. The BAO pure correction ($\Dchi=0.022$) is evaluated at the \LCDM{} best-fit $\Om$ ($0.297$): the pure correction \emph{at the \LCDM{} best fit}. The SNe pure correction of \cref{sec:sne} is evaluated at the bimetric $\Om$ ($0.30435$): the pure correction \emph{at the bimetric closure}. Under the SNe convention the BAO value is $-0.486$; under the BAO convention the SNe values are $+0.4645$ (full covariance, evaluated at its own best-fit $\Om=0.33142$) and $+1.0433$ (diagonal, at its own best-fit $\Om=0.35024$). The qualitative conclusions are independent of the convention; ratios between probes are not. \Cref{fig:mechanism}b shows the decomposition of the per-probe $\Dchi$ in the SNe convention.

Profiling $f=c/(H_0r_d)$ removes the constant part of any rescaling of the background. The same uniform rescaling $E\to AE$ is absorbed by SNe through the absolute magnitude $M$: since $D_L\to D_L/A$, the distance modulus shifts by $\Delta\mu=-5\log_{10}A$, which is a constant absorbed by profiling $M$. Thus BAO and SNe are both insensitive to a constant distance normalization, although the nuisance parameters implementing this degeneracy are different. The SNe penalty does not arise from a lack of such a mechanism. It arises because the anchor-implied effective density differs from the density preferred by Pantheon+, while $M$ cannot remove the redshift-dependent residual of the bimetric background. This is the SNe counterpart of the BAO density-shift effect, not a qualitative distinction in their ability to absorb normalization.

\begin{figure}[t]
\centering
\includegraphics[width=\textwidth]{fig2_mechanism}
\caption{Why BAO is blind at the adopted anchor. (a)~Relative deviation of the exact bimetric $E^2$ from the dark-degeneracy form $E^2_{\rm eff}=\OL+\cos^2\!\theta\,\Om(1+z)^3+\Orad(1+z)^4$ (symmetric-logarithmic axis, sign changes near $z\simeq5$; both signs are shown). Within the BAO and SNe window $z\le2.33$ (shaded) the absolute deviation stays below $3.4\times10^{-5}$ (dashed bounds); it grows to $3.6\times10^{-3}$ at $z=12$. (b)~Decomposition of the per-probe $\Dchi$ against \LCDM{} into the density-mismatch term (orange) and the $E^2(z)$ correction at the bimetric $\Om$ (blue); black bars show the totals. BAO: $+0.665-0.486=0.179$; SNe (full covariance): $2.30+0.64=2.94$; SNe (diagonal): $6.30+1.01=7.32$; CMB: $0.0002$. Generated by \code{src/make\_figures.py}.}\label{fig:mechanism}
\end{figure}
